\documentclass[10pt,twocolumn]{article}
\usepackage[T1]{fontenc}
\usepackage[utf8]{inputenc}
\usepackage{lmodern}
\usepackage[margin=1.8cm,top=2.4cm,bottom=2cm,columnsep=0.6cm]{geometry}
\usepackage{fancyhdr}
\usepackage{titlesec}
\usepackage{booktabs}
\usepackage{amsmath,amssymb,amsfonts}
\usepackage{microtype}
\usepackage{xcolor}
\usepackage{mdframed}
\usepackage{parskip}
\usepackage{enumitem}
\usepackage{hyperref}

\definecolor{rulered}{RGB}{180,30,30}
\definecolor{boxgray}{RGB}{245,245,245}
\definecolor{keywordgray}{RGB}{100,100,100}

\pagestyle{fancy}
\fancyhf{}
\fancyhead[L]{\textsc{\small Claude Mag}}
\fancyhead[C]{\small\textit{ICML 2026 Special Edition: Uncertainty \& Conformal Prediction}}
\fancyhead[R]{\small 2026-07-08}
\fancyfoot[C]{\small\thepage}
\renewcommand{\headrulewidth}{0.8pt}

\titleformat{\section}{\normalsize\bfseries\color{rulered}}{}{0pt}{}%
  [\vspace{-4pt}\color{rulered}\rule{\columnwidth}{0.4pt}]
\titlespacing{\section}{0pt}{8pt}{4pt}

\hypersetup{colorlinks=true,urlcolor=rulered,linkcolor=black}

\begin{document}
\setlength{\parindent}{0pt}
\setlength{\parskip}{4pt}

{\small\color{keywordgray}\textbf{Keywords:}
  conformal prediction \textbullet{}
  decision theory \textbullet{}
  minimax risk \textbullet{}
  coverage guarantees}\par\vspace{4pt}

{\LARGE\bfseries\raggedright
  Optimal Decision-Making Based on Prediction Sets}\par\vspace{4pt}

{\large\itshape\raggedright
  Risk-Optimal Conformal Prediction Derives the Policy and the Set
  from the Same Minimax Objective}\par\vspace{6pt}

{\small\textbf{By} Tao Wang \& Edgar Dobriban
  (University of Pennsylvania, Wharton Statistics and Data Science)}\par\vspace{2pt}

{\small\color{keywordgray}arXiv:2602.00989 \textbullet{}
  ICML 2026 Oral}
\par\vspace{2pt}\noindent{\color{rulered}\rule{\columnwidth}{0.6pt}}\vspace{6pt}

A conformal prediction set guarantees that the true label falls inside it
with probability at least $1-\alpha$---but says nothing about what to
\emph{do} once you have it. This paper derives the minimax-optimal
downstream decision policy for a fixed prediction set, then works backward
to the optimal set construction itself, unifying both into a single
algorithm: \textbf{Risk-Optimal Conformal Prediction (ROCP)}. The result
consistently cuts worst-case risk and critical mistake rates relative to
standard max-min and best-response baselines, with the largest gains where
out-of-set errors are most costly.

\begin{mdframed}[backgroundcolor=boxgray,linecolor=rulered,linewidth=1pt,
    innerleftmargin=6pt,innerrightmargin=6pt,
    innertopmargin=6pt,innerbottommargin=6pt]
\textbf{\small AT A GLANCE}\par\vspace{4pt}
\begin{description}[leftmargin=0pt,labelsep=4pt,itemsep=2pt,parsep=0pt,
    font=\small\bfseries]
  \item[Problem:] Coverage guarantees say nothing about how to act
    optimally on a prediction set for a downstream decision.
  \item[Why it matters:] Prediction sets are increasingly used to drive
    real decisions (triage, risk limits, control); an uncalibrated
    decision rule can waste or misuse the coverage guarantee.
  \item[Method:] Minimax decision theory over distributions consistent
    with the coverage constraint, solved jointly for the optimal policy
    and the optimal prediction set.
  \item[Result:] ROCP reduces worst-case risk and critical mistake rates
    versus max-min and best-response baselines, especially when
    out-of-set errors are costly.
  \item[Evidence:] Comparisons against risk-averse max-min policies and
    naive best-response policies across decision tasks with varying
    error costs.
\end{description}
\end{mdframed}

\section{The Big Picture}

Conformal prediction has become the standard way to wrap any black-box
model with a distribution-free coverage guarantee: the resulting set
$C(x)$ contains the true label with probability at least $1-\alpha$.
But in practice, sets are rarely the end product---they feed a downstream
decision. A clinician acts on a set of possible diagnoses; a control
system acts on a set of possible next states. Until now, the choice of
\emph{how} to act on $C(x)$, and even how $C(x)$ should be shaped for
that purpose, has been left to ad hoc heuristics disconnected from the
coverage guarantee itself.

This paper treats the full pipeline---set construction and downstream
decision---as one decision-theoretic problem, and solves for both pieces
jointly under a minimax risk objective.

\section{Why Existing Approaches Fall Short}

Two heuristics dominate practice. \textbf{Best-response} policies act as
if the point prediction were exact, ignoring the coverage guarantee
outright---leaving decisions exposed exactly when the true label lands
elsewhere in the set. \textbf{Max-min} policies plan for the single
worst label anywhere in $C(x)$, regardless of how likely the coverage
guarantee actually makes it---wasting performance hedging against
possibilities the guarantee doesn't privilege. Neither is derived from
the same objective that produced the set in the first place, so neither
is calibrated to the decision task.

\section{Core Mechanism}

The paper first fixes an arbitrary prediction set $C(x)$ and solves for
the minimax-optimal policy: the one minimizing worst-case expected loss
over all distributions still consistent with the set's coverage
guarantee. That optimal policy balances the worst-case loss for labels
inside $C(x)$ against a penalty for the residual probability mass that
could fall outside it. Holding this policy fixed, the authors then solve
backward for the prediction set that minimizes the resulting robust risk
subject to the coverage constraint---yielding a set shaped for the
decision task, not merely for coverage.

\section{Technical Formulation}

For decision $a \in \mathcal{A}$ and loss $\ell(a,y)$, define the robust
risk of policy $\pi$ given set $C$ as
\[
R(\pi, C) = \sup_{Q:\, Q(y \in C(x)) \ge 1-\alpha}
  \mathbb{E}_Q\big[\ell(\pi(C(x)), y)\big].
\]
The minimax-optimal $\pi^{\ast}$ balances
$\sup_{y \in C(x)} \ell(\pi^{\ast}, y)$ against an $\alpha$-weighted
penalty for loss outside $C(x)$. The optimal set $C^{\ast}(x)$ then
minimizes $R(\pi^{\ast}, C)$ subject to $P(y \in C(x)) \ge 1-\alpha$.
ROCP estimates the population quantities needed for $\pi^{\ast}$ and
$C^{\ast}$ using any black-box probabilistic model, then calibrates a
nonconformity threshold on held-out data so that $C^{\ast}(x)$ retains
exact finite-sample marginal coverage under exchangeability.

\section{Learning or Inference Procedure}

\begin{enumerate}[itemsep=2pt,parsep=0pt]
  \item Fit a black-box probabilistic model to estimate the population
    quantities entering the oracle minimax policy $\pi^{\ast}$.
  \item Construct the candidate set $C^{\ast}(x)$ that minimizes the
    resulting robust risk under the model's estimates.
  \item Calibrate a nonconformity threshold on a held-out calibration
    set to enforce exact finite-sample marginal coverage.
  \item At decision time, apply $\pi^{\ast}$ to the calibrated set
    $C^{\ast}(x)$ to produce the action $a$.
\end{enumerate}

\section{What the Guarantee Says}

Regardless of how well the black-box model estimates the oracle
quantities, the calibration step guarantees
$P(y \in C^{\ast}(x)) \ge 1-\alpha$ in finite samples under
exchangeability alone. The risk-optimality of $\pi^{\ast}$ and
$C^{\ast}$ is with respect to the worst-case distribution consistent
with that guarantee---so the decision quality improves as the
underlying model improves, while the coverage guarantee never depends
on the model being correct.

\section{Experimental Findings}

Across decision tasks with varying costs for out-of-set errors, ROCP
reduces both worst-case risk and the rate of critical (high-cost)
mistakes compared to the max-min baseline (over-conservative) and the
best-response baseline (under-protected). The advantage widens as the
cost of an out-of-set error grows, precisely the regime where a
principled decision policy matters most.

\section{Ablations and Interpretation}

\textbf{Cost sensitivity:} When out-of-set costs are low, ROCP, max-min,
and best-response converge in performance---the decision-theoretic
refinement matters most when mistakes are expensive.

\textbf{Set size:} ROCP's sets are not simply smaller or larger than
standard conformal sets; they are reshaped toward the decision task,
sometimes excluding low-cost-if-missed labels in favor of tighter
coverage where mistakes are expensive.

\textbf{Model quality:} Decision quality scales with the black-box
model's accuracy, while the coverage guarantee remains exact even under
model misspecification---the two guarantees are cleanly separated.

\section*{Reference}

Tao Wang, Edgar Dobriban.
\textbf{Optimal Decision-Making Based on Prediction Sets}.
arXiv:2602.00989, February 2026.
\url{https://arxiv.org/abs/2602.00989}

\end{document}
